\chapter{Các dịch vụ AI}

\section{Dịch vụ khử nhiễu -- Noise Reduction Service}

\subsection{Tổng quan}

Dịch vụ nhận file WAV chứa giọng nói lẫn nhiễu từ ESP32 và trả về audio đã khử nhiễu.
Triển khai trên Flask 2.3+ với TensorFlow 2.13+, kiến trúc layered (API → Service → Model).

\begin{itemize}
\item \textbf{API:} \texttt{POST /denoise} -- multipart/form-data, file WAV tối đa 50 MB.
\item \textbf{Health check:} \texttt{GET /health} -- trạng thái model và thông số.
\item \textbf{Singleton Pattern:} Model được load một lần khi khởi động, cache trong bộ nhớ.
\end{itemize}

\subsection{Mô hình DTLN}

DTLN (Dual-signal Transformation LSTM Network) là mô hình khử nhiễu giọng nói
được đề xuất bởi Westhausen và Meyer (2020), thiết kế cho xử lý thời gian thực
trên thiết bị tài nguyên hạn chế.

\begin{table}[htbp]
\centering
\caption{Các thông số chính của mô hình DTLN}
\label{tab:dtln-params}
\begin{tabular}{>{\raggedright\arraybackslash}p{0.45\linewidth} >{\centering\arraybackslash}p{0.45\linewidth}}
\toprule
\rowcolor{headerblue} \textcolor{white}{\textbf{Thông số}} & \textcolor{white}{\textbf{Giá trị}} \\
\midrule
\rowcolor{rowgray} Tần số lấy mẫu & 16.000 Hz, mono \\
Kích thước khung (blockLen) & 512 mẫu (32 ms) \\
\rowcolor{rowgray} Bước dịch khung (block\_shift) & 128 mẫu (8 ms), overlap 75\% \\
Số bin tần số sau STFT & 257 (= blockLen / 2 + 1) \\
\rowcolor{rowgray} Số đơn vị LSTM mỗi lớp & 128 units \\
Số lớp LSTM mỗi giai đoạn & 2 lớp \\
\rowcolor{rowgray} Dropout & 0,25 \\
Số bộ lọc Encoder Conv1D & 256 \\
\rowcolor{rowgray} Hàm kích hoạt mask & Sigmoid {[}0, 1{]} \\
Hàm mất mát & Negative SNR \\
\rowcolor{rowgray} Tổng số tham số & $\approx$ 1,8 triệu \\
Kích thước file model & 3,9 MB (\texttt{DTLN\_vivos\_best.h5}) \\
\rowcolor{rowgray} Độ trễ xử lý & $\approx$ 32 ms (real-time capable) \\
\bottomrule
\end{tabular}
\end{table}

\subsection{Kiến trúc hai giai đoạn}

\textbf{Giai đoạn 1 -- Miền tần số (STFT Domain):}
Tín hiệu đầu vào được chia thành các khung 512 mẫu, áp dụng STFT để thu được
257 bin tần số. Phổ biên độ $|X(f,t)|$ được chuẩn hóa qua InstantLayerNormalization,
xử lý qua 2 lớp LSTM (128 units) với Dropout 0,25, rồi qua Dense + Sigmoid để tạo
mask $M_1 \in [0,1]$. Mask được nhân với phổ biên độ gốc, kết hợp lại với phổ pha
và biến đổi ngược (iFFT + Overlap-and-Add) để thu được tín hiệu $y_1(t)$.

\begin{table}[htbp]
\centering
\caption{Cấu trúc Separation Kernel -- Giai đoạn 1}
\label{tab:sep-kernel}
\begin{tabular}{>{\raggedright\arraybackslash}p{0.2\linewidth} >{\centering\arraybackslash}p{0.35\linewidth} >{\raggedright\arraybackslash}p{0.35\linewidth}}
\toprule
\rowcolor{headerblue} \textcolor{white}{\textbf{Thành phần}} & \textcolor{white}{\textbf{Thông số}} & \textcolor{white}{\textbf{Vai trò}} \\
\midrule
\rowcolor{rowgray} LSTM lớp 1 & 128 units, return\_seq=True & Đặc trưng tần số cục bộ \\
Dropout & 0,25 & Giảm overfitting \\
\rowcolor{rowgray} LSTM lớp 2 & 128 units, return\_seq=True & Phụ thuộc thời gian dài hạn \\
Dense & 257 đơn vị & Ánh xạ về không gian bin tần số \\
\rowcolor{rowgray} Sigmoid & Đầu ra {[}0, 1{]} & Tạo mask $M_1(f,t)$ \\
\bottomrule
\end{tabular}
\end{table}

\textbf{Giai đoạn 2 -- Miền đặc trưng (Feature Domain):}
Tín hiệu $y_1(t)$ được chiếu sang không gian 256 chiều qua Conv1D point-wise
(kernel=1, stride=1). Sau InstantLayerNormalization và Separation Kernel thứ hai
(cùng cấu trúc LSTM), Dense + Sigmoid tạo mask $M_2$ trong không gian đặc trưng.
Giải mã qua Conv1D point-wise (256→512) với causal padding, rồi Overlap-and-Add
tạo tín hiệu sạch $\hat{y}(t)$.

\begin{table}[htbp]
\centering
\caption{So sánh hai giai đoạn xử lý DTLN}
\label{tab:dual-signal}
\begin{tabular}{>{\raggedright\arraybackslash}p{0.2\linewidth} >{\centering\arraybackslash}p{0.33\linewidth} >{\centering\arraybackslash}p{0.33\linewidth}}
\toprule
\rowcolor{headerblue} \textcolor{white}{\textbf{Tiêu chí}} & \textcolor{white}{\textbf{STFT (GĐ 1)}} & \textcolor{white}{\textbf{Conv1D Encoder (GĐ 2)}} \\
\midrule
\rowcolor{rowgray} Loại biến đổi & Cố định (fixed basis) & Học từ dữ liệu \\
Chiều đầu ra & 257 bin tần số & 256 đặc trưng ẩn \\
\rowcolor{rowgray} Ý nghĩa vật lý & Rõ ràng, diễn giải được & Trừu tượng \\
Vai trò & Loại bỏ nhiễu tổng thể & Tinh chỉnh, phục hồi chi tiết \\
\bottomrule
\end{tabular}
\end{table}

\subsection{Hàm mất mát Negative SNR}

Mô hình được tối ưu hóa bằng hàm mất mát:
\begin{equation}
\text{Loss} = -10 \cdot \log_{10}\left(\frac{E[s^2]}{E[(s - \hat{s})^2 + \epsilon]}\right)
\end{equation}
trong đó $s$ là tín hiệu sạch, $\hat{s}$ là tín hiệu ước lượng, $\epsilon = 10^{-7}$.

\subsection{Huấn luyện trên VIVOS}

\begin{table}[htbp]
\centering
\caption{Cấu hình huấn luyện DTLN trên VIVOS}
\label{tab:training-config}
\begin{tabular}{>{\raggedright\arraybackslash}p{0.3\linewidth} >{\centering\arraybackslash}p{0.2\linewidth} >{\raggedright\arraybackslash}p{0.35\linewidth}}
\toprule
\rowcolor{headerblue} \textcolor{white}{\textbf{Tham số}} & \textcolor{white}{\textbf{Giá trị}} & \textcolor{white}{\textbf{Giải thích}} \\
\midrule
\rowcolor{rowgray} Dữ liệu sạch & VIVOS & Tiếng Việt, $>$ 1 GB \\
Dữ liệu nhiễu & Microsoft DNS & Kết hợp random SNR \\
\rowcolor{rowgray} Batch size & 32 & Cân bằng GPU/gradient \\
Độ dài đoạn & 3 giây & Phù hợp VIVOS (3--5s) \\
\rowcolor{rowgray} Epoch tối đa & 50 & Kết hợp dừng sớm \\
Learning rate & $10^{-3}$ & ReduceLROnPlateau \\
\rowcolor{rowgray} Early Stopping & 10 epoch & Tránh overfitting \\
Optimizer & Adam, clipnorm=3.0 & Tránh bùng nổ gradient \\
\bottomrule
\end{tabular}
\end{table}

\subsection{Pipeline xử lý}

Luồng xử lý khi nhận request:
\begin{enumerate}
\item Đọc file WAV qua SoundFile, validate sample rate (16 kHz).
\item Convert stereo→mono (nếu cần), cast sang float32.
\item Thêm batch dimension, chạy \texttt{model.predict()}.
\item Clip output vào $[-1, 1]$, ghi file WAV đã khử nhiễu.
\item Trả file qua HTTP response, xóa file tạm.
\end{enumerate}

\section{Dịch vụ nhận dạng giọng nói -- Voice Rephraze Service}

\subsection{Tổng quan}

Dịch vụ chuyển đổi giọng nói thành văn bản và sinh mô tả sản phẩm theo
phong cách viết tùy chọn. Triển khai trên FastAPI với AsyncOpenAI client.

\subsection{Speech-to-Text}

\begin{itemize}
\item \textbf{Mô hình:} OpenAI gpt-4o-transcribe (cloud API).
\item \textbf{Định dạng hỗ trợ:} WAV, MP3, FLAC, OGG, WebM, M4A.
\item \textbf{Ngôn ngữ:} Tiếng Việt (đa vùng miền).
\item \textbf{API:} \texttt{POST /stt} -- upload file audio, trả về text.
\item \textbf{Lịch sử phát triển:} PhoWhisper-tiny (local) → whisper-1 (OpenAI) → gpt-4o-transcribe.
\end{itemize}

\subsection{Sinh mô tả sản phẩm}

\begin{itemize}
\item \textbf{Mô hình:} GPT-4o-mini (temperature 0.75, max 600 tokens).
\item \textbf{API:} \texttt{POST /gen} -- nhận \texttt{product\_description} (text từ STT)
  và \texttt{style} (phong cách viết).
\item \textbf{Prompt engineering:} System prompt hướng dẫn AI đóng vai copywriter
  thương mại điện tử, trích xuất thông tin sản phẩm từ text lộn xộn của STT,
  bỏ qua từ lấp, lặp lại, lỗi phát âm; viết 3--5 câu ngắn gọn theo phong cách yêu cầu.
\end{itemize}

\begin{table}[htbp]
\centering
\caption{Các phong cách viết mô tả sản phẩm}
\label{tab:writing-styles}
\begin{tabular}{>{\raggedright\arraybackslash}p{0.25\linewidth} >{\raggedright\arraybackslash}p{0.6\linewidth}}
\toprule
\rowcolor{headerblue} \textcolor{white}{\textbf{Phong cách}} & \textcolor{white}{\textbf{Mô tả}} \\
\midrule
\rowcolor{rowgray} Văn minh & Lịch sự, trang trọng, ngôn ngữ chuẩn mực \\
Chuyên nghiệp & Kỹ thuật, rõ ràng, tập trung thông số \\
\rowcolor{rowgray} Thân thiện & Ấm áp, gần gũi, dễ tiếp cận \\
Văn hóa người Tày & Yếu tố văn hóa dân tộc, hình ảnh đặc trưng \\
\bottomrule
\end{tabular}
\end{table}

\subsection{Lưu trữ và quản lý}

Các phong cách viết được lưu trong SQLite (\texttt{styles.db}) với schema đơn giản
(id, name, description). Người dùng có thể thêm phong cách mới qua
\texttt{POST /add-style}.

\section{Dịch vụ phân loại sản phẩm -- Product Classification Service}

\subsection{Tổng quan}

Dịch vụ phân loại text mô tả sản phẩm vào 32 danh mục thương mại điện tử.
Triển khai trên FastAPI 0.115 với PyTorch 2.5.1 và HuggingFace Transformers 4.46.3.

\subsection{Mô hình XLM-RoBERTa}

\begin{itemize}
\item \textbf{Model:} \texttt{Lezh1n/xlm-roberta-ecommerce-classifier} (HuggingFace Hub).
\item \textbf{Base:} XLM-RoBERTa-base (Facebook AI) -- hỗ trợ đa ngôn ngữ.
\item \textbf{Task:} Multi-class sequence classification (32 classes).
\item \textbf{Kích thước:} $\approx$ 1,1 GB, max sequence 512 tokens.
\item \textbf{Accuracy:} 90,14\%. \textbf{F1-score:} 90,00\%.
\end{itemize}

\subsection{API Endpoints}

\begin{table}[htbp]
\centering
\caption{API endpoints dịch vụ phân loại}
\label{tab:classify-api}
\begin{tabular}{>{\raggedright\arraybackslash}p{0.35\linewidth} >{\raggedright\arraybackslash}p{0.55\linewidth}}
\toprule
\rowcolor{headerblue} \textcolor{white}{\textbf{Endpoint}} & \textcolor{white}{\textbf{Chức năng}} \\
\midrule
\rowcolor{rowgray} \texttt{POST /classify} & Phân loại 1 sản phẩm, trả top-K predictions \\
\texttt{POST /classify/batch} & Batch tối đa 100 items (batch\_size=32) \\
\rowcolor{rowgray} \texttt{GET /categories} & Danh sách 32 danh mục \\
\texttt{GET /model-info} & Metadata và metrics mô hình \\
\rowcolor{rowgray} \texttt{GET /health} & Health check \\
\bottomrule
\end{tabular}
\end{table}

\subsection{32 danh mục sản phẩm}

Electronics, Computers \& Networking, Mobile Phones \& Tablets, Fashion \& Clothing,
Shoes \& Footwear, Bags \& Luggage, Watches, Jewelry, Fashion Accessories,
Beauty \& Personal Care, Health \& Wellness, Sports \& Outdoors, Home Furniture,
Home Decor \& Lighting, Kitchen \& Dining, Bedding \& Bath, Large Appliances,
Small Appliances, Grocery \& Food, Baby \& Maternity, Toys \& Games, Books \& Media,
Stationery \& Office Supplies, Pet Supplies, Automotive \& Motorcycle,
Tools \& Hardware, Garden \& Outdoor Living, Musical Instruments,
Video Games \& Gaming, Software \& Digital Goods, Arts \& Crafts,
Industrial \& Commercial.

\subsection{Hiệu năng}

\begin{table}[htbp]
\centering
\caption{Benchmark hiệu năng phân loại sản phẩm}
\label{tab:classify-benchmark}
\begin{tabular}{>{\raggedright\arraybackslash}p{0.35\linewidth} >{\centering\arraybackslash}p{0.2\linewidth} >{\centering\arraybackslash}p{0.2\linewidth}}
\toprule
\rowcolor{headerblue} \textcolor{white}{\textbf{Scenario}} & \textcolor{white}{\textbf{Latency}} & \textcolor{white}{\textbf{Throughput}} \\
\midrule
\rowcolor{rowgray} Single inference (CPU) & $\approx$ 45 ms & 22 req/s \\
Batch 10 items (CPU) & $\approx$ 180 ms & 55 items/s \\
\rowcolor{rowgray} Single inference (GPU) & $\approx$ 15 ms & 66 req/s \\
Batch 100 items (GPU) & $\approx$ 800 ms & 125 items/s \\
\bottomrule
\end{tabular}
\end{table}

\subsection{Pipeline suy luận}

\begin{enumerate}
\item Client gửi text mô tả sản phẩm qua POST request.
\item Validate input bằng Pydantic schema.
\item Tokenize input (max 512 tokens) qua AutoTokenizer.
\item Forward pass sinh logits cho 32 danh mục.
\item Softmax → top-K predictions với category ID, name, display name, score.
\item Trả response kèm inference time.
\end{enumerate}

Model được load khi khởi động qua FastAPI Lifespan Context Manager,
tự động detect GPU (fallback CPU), set evaluation mode.
